\subtitle{\bf Aula 10 - Motivação}
\begin{frame}[t,plain]
    \titlepage
\end{frame}

\begin{frame}{Objetivos}

    \begin{itemize}
        \item
        Apresentar a motivação para os algoritmos LR(1) e LALR.
        \item
        Apresentar o algoritmo LR(1).
        \item
        Apresentar o algoritmo LALR como uma otimização de LR(1).
    \end{itemize}
\end{frame}


\section{Motivação: Além do SLR}\label{motivauxe7uxe3o-aluxe9m-do-slr}

\begin{frame}{Limitações}

    \begin{itemize}

        \item
        \textbf{LR(0)}: insere reduce em todos os terminais

        \begin{itemize}

            \item
            Muitos conflitos
        \end{itemize}
        \item
        \textbf{SLR}: usa o follow(A)

        \begin{itemize}

            \item
            Ainda impreciso
        \end{itemize}
    \end{itemize}
\end{frame}

\begin{frame}{Limitações}

    \textbf{Exemplo problemático} com recursão à direita:

    \[\begin{array}{lcl}
    E &\to& T + E \mid T \\
    T &\to& n \mid (E)
\end{array}\]
\end{frame}

\begin{frame}{Limitações}

    Estado após reconhecer \(T\) no estado inicial:

    \[\begin{array}{l}
    E \to T . + E \\
    E \to T .
\end{array}\]
\end{frame}

\begin{frame}[fragile]{O Conflito Shift-Reduce}

    \begin{itemize}

        \item
        Item \(E \to T .\) está completo

        \begin{itemize}

            \item
            LR(0) insere reduce \(E \to T\) em \textbf{todos} os terminais.
        \end{itemize}
        \item
        Item \(E \to T . + E\) exige \textbf{shift} em
        \texttt{+}

        \begin{itemize}

            \item
            Resultado: \textbf{conflito shift-reduce} em
            \texttt{+}
        \end{itemize}
    \end{itemize}
\end{frame}

\begin{frame}[fragile]{O Conflito Shift-Reduce}

    \begin{itemize}

        \item
        A raiz do problema

        \begin{itemize}

            \item
            Quando o próximo token é \texttt{+}, ainda pode vir
            \(+E\), então \ldots{}
            \item
            \textbf{Não} devemos reduzir.
            \item
            Quando é \texttt{$} ou \texttt{)},
            devemos reduzir.
        \end{itemize}
    \end{itemize}
\end{frame}

\begin{frame}{O Conflito de Shift-Reduce}

    \begin{itemize}

        \item
        Para tomar essa decisão, precisamos analisar o próximo token de
        entrada.

        \begin{itemize}

            \item
            Lookahead.
        \end{itemize}
        \item
        Um único token de lookahead local resolve o conflito

        \begin{itemize}

            \item
            Essa é a ideia do \textbf{LR(1)}.
        \end{itemize}
    \end{itemize}
\end{frame}



\begin{frame}{Itens LR(1)}

    Um item LR(1) \([A \to \beta . \gamma, a]\) carrega:

    \begin{itemize}

        \item
        \(A \to \beta\gamma\): produção sendo reconhecida
        \item
        \(.\) (ponto): o que já foi analisado
        \item
        \(a\): \textbf{lookahead} --- terminal que pode aparecer
        \textbf{imediatamente após} \(A\) neste contexto específico
    \end{itemize}
\end{frame}

\begin{frame}{Redução}

    \begin{itemize}

        \item
        Quando o item está completo (\(\gamma = \lambda\)):

        \begin{itemize}

            \item
            reduzir \(A \to \beta\) \textbf{somente} na coluna \(a\).
        \end{itemize}
    \end{itemize}
\end{frame}

\begin{frame}{Redução}

    \begin{itemize}

        \item
        Lookahead é \textbf{calculado localmente} para cada item, não
        globalmente como no SLR.
    \end{itemize}
\end{frame}

\begin{frame}{Fechamento LR(1)}

    \begin{itemize}

        \item
        Para item \([A \to \alpha . B \beta, a]\), adicionar
        \([B \to . \gamma, b]\) para toda produção \(B \to \gamma\) e todo
        \(b \in \mathrm{FIRST}(\beta\, a)\).
    \end{itemize}
\end{frame}

\begin{frame}{Fechamento LR(1)}

    \begin{itemize}

        \item
        Se \(\beta\) não é anulável: \(b \in \mathrm{FIRST}(\beta)\)
        \item
        Se \(\beta\) é anulável ou vazio: \(b\) inclui também o lookahead
        \(a\)
    \end{itemize}
\end{frame}

\begin{frame}{Construção dos Estados LR(1)}

    \begin{itemize}
        \item
        Estado inicial: fechamento de \(\{[S' \to . S, \$]\}\)
        \item
        O lookahead é transferido \textbf{sem alteração} --- Ele depende do
        contexto em que o item foi criado.
    \end{itemize}
\end{frame}

\begin{frame}[fragile]{Tabela LR(1)}

    \begin{itemize}

        \item
        \textbf{Shift}: \([A \to \alpha . a \beta, b]\) com \(a\) terminal ->
        \texttt{action[I, a] = Shift j}
    \end{itemize}
\end{frame}

\begin{frame}[fragile]{Tabela LR(1)}

    \begin{itemize}

        \item
        \textbf{Reduce}: \([A \to \alpha ., a]\) completo ->
        \texttt{action[I, a] = Reduce(A -> α)} \textbf{apenas
        na coluna \(a\)}
    \end{itemize}
\end{frame}

\begin{frame}[fragile]{Tabela LR(1)}

    \begin{itemize}

        \item
        \textbf{Accept}: \([S' \to S ., \$]\) ->
        \texttt{action[I, $] = Accept}
    \end{itemize}
\end{frame}

\begin{frame}{Tabela LR(1)}

    \begin{itemize}

        \item
        Cada reduce aparece \textbf{apenas na coluna do lookahead do item} ---
        não em todas.
    \end{itemize}
\end{frame}


\section{Exemplo}\label{exemplo}

\begin{frame}{Gramática}

    \[\begin{array}{lcl}
    E &\to& E + T \mid T \\
    T &\to& T * F \mid F \\
    F &\to& (E) \mid n
\end{array}\]
\end{frame}

\begin{frame}{Estado 0}

    \begin{itemize}

        \item
        Fechamento de \(\{[S' \to . E, \$]\}\):
    \end{itemize}

    \[\begin{array}{ll}
    S' \to . E & \$ \\
    E \to . E + T & \$, + \\
    E \to . T & \$, + \\
    T \to . T * F & \$, +, * \\
    T \to . F & \$, +, * \\
    F \to . (E) & \$, +, *, ) \\
    F \to . n & \$, +, *, )
\end{array}\]
\end{frame}

\begin{frame}{Estado 1}

    --- \(\text{goto}(I_0, E)\):

    \[\begin{array}{ll}
    S' \to E . & \$ \\
    E \to E . + T & \$, +
\end{array}\]
\end{frame}

\begin{frame}{Estado 2}

    --- \(\text{goto}(I_0, T)\):

    \[\begin{array}{ll}
    E \to T . & \$, + \\
    T \to T . * F & \$, +, *
\end{array}\]
\end{frame}

\begin{frame}{Estado 3}

    --- \(\text{goto}(I_0, F)\):

    \[\begin{array}{ll}
    T \to F . & \$, +, *
\end{array}\]
\end{frame}

\begin{frame}{Tabela LR(1)}

    \begin{table}
        \begin{tabular}[]{@{}llllllllll@{}}
            \toprule
            Estado & \(n\) & \(+\) & \(*\) & \((\) & \()\) & \(\$\) & \(E\) & \(T\)
            & \(F\) \\
            \midrule

            0 & s5 & & & s4 & & & 1 & 2 & 3 \\
            1 & & s6 & & & & acc & & & \\
            2 & & r E->T & s7 & & & r E->T & & & \\
            3 & & r T->F & r T->F & & & r T->F & & & \\
            4 & s5 & & & s4 & & & 8 & 9 & 3 \\
            5 & & r F->n & r F->n & & r F->n & r F->n & & & \\
            \bottomrule
        \end{tabular}
    \end{table}
\end{frame}

\begin{frame}{Tabela LR(1)}

    Cada ação de reduce aparece \textbf{apenas nas colunas do lookahead do
    item} --- não em todas.
\end{frame}


\section{O Algoritmo LALR(1)}\label{o-algoritmo-lalr1}

\begin{frame}{Motivação}

    \begin{itemize}

        \item
        LR(1) gera estados distintos para o mesmo conjunto de itens LR(0)
        quando os lookaheads diferem.
    \end{itemize}
\end{frame}

\begin{frame}{Motivação}

    \begin{itemize}

        \item
        \textbf{Exemplo:} os estados 3 e 10 da gramática de expressões têm o
        mesmo núcleo \(T \to F .\):

        \begin{itemize}

            \item
            Estado 3 (fora de parênteses): lookaheads \(\{\$, +, *\}\)
            \item
            Estado 10 (dentro de parênteses): lookaheads \(\{\$, +, *, )\}\)
        \end{itemize}
    \end{itemize}
\end{frame}

\begin{frame}{Motivação}

    \begin{itemize}

        \item
        Para gramáticas reais, o número de estados LR(1) pode ser uma
        \textbf{ordem de grandeza maior} que LR(0).
    \end{itemize}
\end{frame}

\begin{frame}{A Ideia LALR}

    \begin{itemize}

        \item
        \textbf{Observação}: dois estados com o mesmo \textbf{núcleo} têm a
        mesma estrutura de shifts e gotos.
    \end{itemize}
\end{frame}

\begin{frame}{A Ideia LALR}

    \begin{itemize}
        \item
        \textbf{Solução}: fundir todos os estados LR(1) com o mesmo núcleo,
        \textbf{unindo} seus lookaheads.
        \item
        \textbf{Resultado}: \(|\text{estados LALR}| = |\text{estados LR(0)}|\)
        mas lookaheads mais precisos que SLR.
    \end{itemize}
\end{frame}

\begin{frame}{Algoritmo LALR}

    \begin{itemize}

        \item
        Dois passos:
    \end{itemize}

    \begin{enumerate}

        \item
        Construir todos os estados LR(1)
        \item
        Fundir estados com o mesmo núcleo, realizando a união de seus
        lookaheads.
    \end{enumerate}
\end{frame}

\begin{frame}{Problemas}

    \begin{itemize}
        \item
        A fusão pode introduzir conflitos \textbf{reduce-reduce}:
        \item
        Se dois estados tinham lookaheads disjuntos, a união pode fazer dois
        itens completos reivindicar a mesma coluna.
    \end{itemize}
\end{frame}

\begin{frame}{Problemas}

    \begin{itemize}

        \item
        Quando isso ocorre: gramática é LR(1) mas \textbf{não LALR(1)}
        \item
        Conflitos \textbf{shift-reduce} nunca são introduzidos pela fusão
        (shifts dependem do núcleo, não do lookahead)
    \end{itemize}
\end{frame}

\begin{frame}{Problemas}

    \begin{itemize}

        \item
        Na prática: a imensa maioria das gramáticas de linguagens de
        programação é LALR(1).
    \end{itemize}
\end{frame}


\section{Exemplo: Tabela LALR(1)}\label{exemplo-tabela-lalr1}

\begin{frame}{Identificação dos Pares com o Mesmo Núcleo}

    \begin{table}
        \begin{tabular}[]{@{}lll@{}}
            \toprule
            Núcleo & Estados LR(1) & Lookaheads unidos \\
            \midrule

            \(T \to F .\) & 3 e 10 & \(\{\$, +, *, )\}\) \\
            \(E \to E + . T\), \ldots{} & 6 e 14 & \(\{\$, +, )\}\) \\
            \(T \to T * . F\), \ldots{} & 7 e 15 & \(\{\$, +, *, )\}\) \\
            \(E \to E + T .\), \(T \to T . * F\) & 11 e 16 & \(\{\$, +, )\}\) \\
            \(T \to T * F .\) & 12 e 17 & \(\{\$, +, *, )\}\) \\
            \bottomrule
        \end{tabular}
    \end{table}

    Denominamos: \textbf{A} = 3+10, \textbf{B} = 6+14, \textbf{C} = 7+15,
    \textbf{D} = 11+16, \textbf{E} = 12+17.
\end{frame}

\begin{frame}{Tabela LALR(1) para Expressões}

    \begin{table}
        \begin{tabular}[]{@{}
            >{\raggedright\arraybackslash}p{(\columnwidth - 18\tabcolsep) * \real{0.1224}}
            >{\raggedright\arraybackslash}p{(\columnwidth - 18\tabcolsep) * \real{0.0612}}
            >{\raggedright\arraybackslash}p{(\columnwidth - 18\tabcolsep) * \real{0.1429}}
            >{\raggedright\arraybackslash}p{(\columnwidth - 18\tabcolsep) * \real{0.1429}}
            >{\raggedright\arraybackslash}p{(\columnwidth - 18\tabcolsep) * \real{0.0612}}
            >{\raggedright\arraybackslash}p{(\columnwidth - 18\tabcolsep) * \real{0.1429}}
            >{\raggedright\arraybackslash}p{(\columnwidth - 18\tabcolsep) * \real{0.1429}}
            >{\raggedright\arraybackslash}p{(\columnwidth - 18\tabcolsep) * \real{0.0612}}
            >{\raggedright\arraybackslash}p{(\columnwidth - 18\tabcolsep) * \real{0.0612}}
            >{\raggedright\arraybackslash}p{(\columnwidth - 18\tabcolsep) * \real{0.0612}}@{}}
            \toprule
            \begin{minipage}[b]{\linewidth}\raggedright
                Estado
            \end{minipage} & \begin{minipage}[b]{\linewidth}\raggedright
            \(n\)
        \end{minipage} & \begin{minipage}[b]{\linewidth}\raggedright
        \(+\)
    \end{minipage} & \begin{minipage}[b]{\linewidth}\raggedright
    \(*\)
\end{minipage} & \begin{minipage}[b]{\linewidth}\raggedright
\((\)
\end{minipage} & \begin{minipage}[b]{\linewidth}\raggedright
\()\)
\end{minipage} & \begin{minipage}[b]{\linewidth}\raggedright
\(\$\)
\end{minipage} & \begin{minipage}[b]{\linewidth}\raggedright
\(E\)
\end{minipage} & \begin{minipage}[b]{\linewidth}\raggedright
\(T\)
\end{minipage} & \begin{minipage}[b]{\linewidth}\raggedright
\(F\)
\end{minipage} \\
\midrule

0 & s5 & & & s4 & & & 1 & 2 & A \\
1 & & sB & & & & acc & & & \\
2 & & r E->T & sC & & & r E->T & & & \\
A & & r T->F & r T->F & & r T->F & r T->F & & & \\
4 & s5 & & & s4 & & & 8 & 9 & A \\
5 & & r F->n & r F->n & & r F->n & r F->n & & & \\
B & s5 & & & s4 & & & & D & A \\
C & s5 & & & s4 & & & & & E \\
8 & & sB & & & s13 & & & & \\
D & & r E->E+T & sC & & r E->E+T & r E->E+T & & & \\
E & & r T->T*F & r T->T*F & & r T->T*F & r T->T*F & & & \\
\bottomrule
\end{tabular}
\end{table}
\end{frame}


\section{Conclusão}\label{conclusuxe3o}

\begin{frame}{Conclusão}

    \begin{itemize}
        \item
        Neste capítulo apresentamos os algoritmos LR(1) e LALR.
        \item
        O gerador de analisadores sintáticos Happy utiliza o algoritmo LALR.
    \end{itemize}
\end{frame}
